\chapter{Security}
\label{ch:security}

Security depends on which connection is being considered. A FIX member connects
to a gateway, the gateway authenticates the COMPID with the authentication
service, and an operator manages COMPIDs through a web application. An operator
also halts trading, lifts a halt, moves the venue from one phase of the
trading day to the next, and replaces the day's schedule, by sending requests
to the sequencers over a connection of its own. These links do not all have the same protection. In particular, SCRAM keeps a COMPID
password out of the gateway-to-authentication-service exchange, but it does not
encrypt a member's connection to the gateway. TLS is available for FIX sessions
but is optional, and the COMPID editor has its own administrator login and
first-run setup.

\section{COMPID authentication}
\label{sec:security:compid_authentication}

A member sends its COMPID and password in the FIX Logon message. The password
is in tag 554 and reaches the FIX gateway. The project calls the gateway's
challenge-response authentication SCRAM-SHA-256. It uses PBKDF2-HMAC-SHA-256,
HMAC-SHA-256, and SHA-256. RFC 7677 specifies the standardized SCRAM-SHA-256
and SCRAM-SHA-256-PLUS SASL mechanisms \cite{rfc7677}. This venue carries its
authentication exchange in project-specific PDU messages, not the RFC's SASL
wire protocol.

The gateway sends the COMPID and a client nonce. The authentication service
returns a server nonce, a salt, and an iteration count. The gateway derives a
client proof from the password and clears the temporary password value
immediately after deriving the salted password. It then sends the proof.
The authentication service checks the proof against the stored
credential and returns a result with a server signature. The gateway verifies
that signature before accepting a successful result. The authentication
service stores derived keys, a salt, and the iteration count, not the
plaintext password.

The connection between the gateway and authentication service is plain TCP,
not TLS. The password is not sent in plaintext on this link: it remains at the
gateway and is used to calculate the SCRAM proof. This protection is specific
to the gateway-to-authentication-service exchange. It does not protect tag 554
on a plain member-to-gateway FIX connection.

\begin{requirement}{R-0155}{The gateway does not send COMPID passwords in plaintext to the authentication service}
  The gateway-to-authentication-service authentication exchange shall not
  transmit a COMPID password in plaintext.
  \rationale{The gateway must authenticate members without exposing their passwords on the network link to the authentication service.}
  \uncovered{no test verifies that a COMPID password is absent from the gateway-to-authentication-service wire messages}
\end{requirement}

\section{Optional TLS for FIX sessions}
\label{sec:security:fix_tls}

The FIX gateway always registers a plain TCP listener. When TLS is enabled in
the gateway configuration, it also registers a TLS listener on a separate
port. A member can therefore connect through TLS or use the plain listener;
members are not required to use TLS.

Some members classify their venue connection as secure under their own
controls and choose plain TCP to avoid the performance cost of encryption.
That is the member's assessment of its connection and risk. It does not mean
the FIX traffic is encrypted by the venue. On a plain connection, tag 554 and
the rest of the FIX session cross the network without TLS protection. A
deployment that enables TLS currently uses TLS 1.2.

\begin{requirement}{R-0154}{Members can choose TLS or plain TCP for FIX sessions}
  A deployment shall be able to offer a TLS listener for FIX sessions while
  retaining the plain TCP listener, so that each member can choose which
  listener to use.
  \rationale{Members have different security controls and performance
  requirements. The gateway must support encrypted FIX sessions without
  requiring every member to use TLS.}
  \uncovered{no test verifies that the FIX gateway offers both listeners when TLS is enabled}
\end{requirement}

\section{Authentication service administration}
\label{sec:security:authentication_service_admin}

The admin service sends credential changes to the authentication service over
its TLS admin connection. Client-certificate verification can be enabled, but
it is disabled by default. The gateway-facing authentication connection is a
separate plain TCP listener and does not use TLS. SCRAM protects the COMPID
password on that connection, but does not encrypt the rest of its traffic.

\section{The operator's connection to the sequencers}
\label{sec:security:operator_connection}

An operator halts trading, lifts a halt, asks for the venue to enter the
next phase of the trading day, and replaces the times that remain in the day's
schedule. Each of these is a request sent to both sequencers, and the leading
sequencer acts on it. Two tools send them:
\texttt{scripts/operator\_halt.py} for a halt and its lift, and
\texttt{scripts/trading\_phase.py} for a phase and the schedule. What the two share, including
the connection and its security, is in \texttt{scripts/operator\_requests.py}.
The procedures that start the venue for development and for testing send the
same requests through the same code.

Each sequencer has a listener for operators, configured in its
\texttt{[operator]} section, and separate from every other listener it has. The
listener accepts only TLS. It requires the client to present a certificate, and
accepts the connection only if that certificate was signed by the operator's
certificate authority, whose certificate is named by
\texttt{tls\_ca\_path}. The sequencer refuses to start if
\texttt{tls\_ca\_path} is empty, because without an authority any client could
halt the venue. A client that presents no certificate, or one signed by another
authority, fails the TLS handshake, and the sequencer reads nothing from it. The
listener uses TLS 1.2 only, with two cipher suites, ECDHE-RSA-AES256-GCM-SHA384
and ECDHE-ECDSA-AES256-GCM-SHA384, as every TLS listener in the framework does.

The tools check the sequencer in turn. They trust only the operator's
certificate authority, not the certificate authorities installed on the machine,
and they check that the sequencer's certificate names the host they connected
to. A program listening in a sequencer's place therefore cannot answer a request
unless it holds a certificate the operator's authority signed for that host.

The deployment script, \texttt{deploy.py}, makes the certificates when the
environment file has an \texttt{[operator\_certificates]} section:

\begin{itemize}
  \item a self-signed certificate authority, \texttt{operator\_ca.crt}, with its
        private key \texttt{operator\_ca.key};
  \item one operator's certificate, \texttt{operator.crt}, with its private key
        \texttt{operator.key}, signed by the authority and usable only to
        identify a client;
  \item for each sequencer, a certificate for its listener, signed by the same
        authority, usable only to identify a server, and naming
        \texttt{localhost}, \texttt{127.0.0.1} and the host the listener binds
        to, or the machine's fully qualified name when it binds to every
        interface; and a copy of the authority's certificate, which the
        sequencer uses to check the operator's certificate.
\end{itemize}

The authority, the operator's certificate and their keys are kept in the
directory the section names, \texttt{etc/operator} under the install directory.
Each private key can be read only by the account that ran the deployment. Every
certificate is valid for ten years. Nothing that exists is replaced unless the
deployment is run with \texttt{--force-certs}, which makes a new authority and
replaces every certificate it signs; each sequencer reads its certificates when
it starts.

The protection has these limits, which a deployment has to take into account:

\begin{itemize}
  \item \textbf{One certificate identifies every operator.} The deployment makes
        a single operator's certificate, so the sequencers cannot tell one
        operator from another. Each request carries the operator's name, which
        each sequencer writes in its Quill log at Warning with the request and
        its answer. The name is what the operator typed, as
        \texttt{--operator}; it is not checked against the certificate, and so
        it records who said they acted, not who did.
  \item \textbf{The certificate decides everything.} A client holding the
        operator's certificate and key may make every request: halt, lift, move
        the phase, and replace the schedule. There is no finer permission.
  \item \textbf{The authority's private key is in the install directory.} Any
        account that can read \texttt{operator\_ca.key} can sign a certificate
        the sequencers accept. It is protected only by the file's permissions.
  \item \textbf{A certificate cannot be withdrawn on its own.} There is no
        certificate revocation list. A lost operator key is dealt with by making
        a new authority with \texttt{--force-certs}, which replaces every
        certificate, and restarting the sequencers.
  \item \textbf{The listener's reach is a deployment setting.} The production
        environment binds the listener to every interface, so which machines
        can reach it is decided by the network, not by the sequencer.
\end{itemize}

\begin{requirement}{R-0176}{Only a client with the operator's certificate can make an operator's request}
  A sequencer shall act on a request to halt trading, to lift a halt, to
  move the venue to another phase of the trading day, or to replace the day's
  schedule only when it arrives over
  TLS from a client that presented a certificate signed by the operator's
  certificate authority. A sequencer configured without that authority shall
  refuse to start.
  \rationale{Each of these requests stops or starts trading for every member.
  A listener that accepted any client would let anyone who can reach it halt
  the venue, or close it for the day. Refusing to start without the authority
  means the protection cannot be left out by mistake.}
  \covers{ha_test:112}
\end{requirement}

\section{COMPID password lifecycle}
\label{sec:security:compid_password_lifecycle}

The COMPID database stores the derived SCRAM keys, salt, and iteration count.
It also has fields for the time a password changed and whether a password
change should be required. The editor can set the change-required flag, and a
password reset updates the change time. The member authentication flow does
not enforce password expiry or the change-required flag: the authentication
service reports that no password change is required. Password expiry and
forced change during FIX logon are therefore not implemented.

\section{COMPID editor and administrator access}
\label{sec:security:compid_editor}

The COMPID editor controls member credentials, account status, gateway
permissions, and session limits. Its administrator accounts are stored
separately from member COMPIDs. Administrator passwords are stored as bcrypt
hashes. An authenticated \texttt{ADMIN} can make changes; a \texttt{VIEWER}
can view data but cannot submit changes. The application code starts a web
server without configuring TLS. Browser-to-editor traffic is therefore not
encrypted by the application itself; deployments need to account for this
when deciding which networks may reach the editor.

When the administrator-user store is empty, the setup page is public and the
first account created receives the \texttt{ADMIN} role. To create it, setup
asks for a one-time secret as well as the username and password. The editor
writes the secret, when it starts with no administrator, to a file beside the
administrator-user store that only the account running the editor can read,
and deletes it when the first account is created.

\begin{requirement}{R-0153}{First administrator setup requires proof of installation access}
  Before the first administrator account is created, setup shall require a
  one-time secret generated for this installation and available only to an
  operator with access to the host. Creating the first account shall consume
  the secret and close first-run setup.
  \rationale{If reaching the setup page is enough to create the first administrator, any
  unauthenticated client that can reach the editor may claim control of member
  credentials and permissions before the operator does. Requiring proof available
  only from the installed host reserves initial administration for the operator.}
  \uncovered{the COMPID editor's unit tests check that setup refuses a wrong or missing secret and that creating the first account deletes it; no scenario starts the editor and tries setup}
\end{requirement}
